\subsection{Prosjektplanlegging}
\label{sec:prosjektplan}
\subsubsection{Arbeidsform og dokumentasjon}
Arbeidet gjennomføres i korte utviklingsrunder med sporbar sammenheng mellom krav, designvalg og tester. CAD-versjoner, bilder, produksjonsinnstillinger, testdata og beslutninger lagres fortløpende. Rapporten oppdateres etter hver hovedfase.

\subsubsection{Fremdriftsplan}
Prosjektet gjennomføres fra uke 39 til innlevering mandag 16. november 2026 (uke 47). Fremdriftsplanen er vist som Gantt-diagram i figur~\ref{fig:gantt}, og aktivitetene er:
\begin{enumerate}
    \item \textbf{Utrede og forstå teorien bak KERS (uke 39).} Sette seg inn i hvordan kinetisk energi kan lagres og gjenvinnes, og etablere det teoretiske grunnlaget for konseptarbeidet.
    \item \textbf{Prosjektstyring og risikovurdering (uke 40).} Planlegge gjennomføringen av prosjektet, fordele arbeid og vurdere risiko med tilhørende tiltak.
    \item \textbf{Konseptutvikling og innledende design (uke 41--43).} Utarbeide produktkrav, utvikle og vurdere konsepter, velge konsept og utforme et første design som kan testes.
    \item \textbf{Testing (uke 43--44).} Fastsette testprosedyre og teste det valgte konseptet. Testingen starter i siste uke av designfasen, slik at funn kan tas med i det videre designarbeidet.
    \item \textbf{Iterering på konsept (uke 44).} Gjennomføre avgrensede forbedringer basert på testresultatene.
    \item \textbf{Løpende rapportskriving (uke 39 til innlevering).} Rapporten skrives parallelt med de øvrige aktivitetene gjennom hele prosjektet.
    \item \textbf{Sluttføring av rapport (uke 45 til innlevering).} Ferdigstille resultater, diskusjon og konklusjoner fra testingen frem mot innleveringen 16. november.
\end{enumerate}
Aktivitetene bygger på hverandre: testing forutsetter et testbart design, og iterasjonen bygger på resultatene fra testingen.

\begin{figure}[htbp]
    \centering
    \begin{ganttchart}[
        x unit = 1.0cm,
        y unit chart = 0.75cm,
        vgrid,
        hgrid,
        title height = 1,
        title label font = \footnotesize,
        bar label font = \footnotesize,
        bar label node/.append style = {align = right, text width = 4.6cm},
        bar/.append style = {fill = NavyBlue!70, draw = NavyBlue},
        bar height = 0.5,
        milestone label font = \footnotesize\bfseries,
        milestone/.append style = {fill = BrickRed, draw = BrickRed},
        milestone label node/.append style = {align = right, text width = 4.6cm},
    ]{1}{9}
        \gantttitle{Uke 39}{1} \gantttitle{Uke 40}{1} \gantttitle{Uke 41}{1}
        \gantttitle{Uke 42}{1} \gantttitle{Uke 43}{1} \gantttitle{Uke 44}{1}
        \gantttitle{Uke 45}{1} \gantttitle{Uke 46}{1} \gantttitle{Uke 47}{1} \\
        \ganttbar{Utrede og forstå teori bak KERS}{1}{1} \\
        \ganttbar{Prosjektstyring og\\risikovurdering}{2}{2} \\
        \ganttbar{Konseptutvikling og\\innledende design}{3}{5} \\
        \ganttbar{Testing}{5}{6} \\
        \ganttbar{Iterering på konsept}{6}{6} \\
        \ganttbar[bar/.append style = {fill = NavyBlue!20}]{Løpende rapportskriving}{1}{8} \\
        \ganttbar{Sluttføring av rapport med\\konklusjoner fra testing}{7}{8} \\
        \ganttmilestone{Innlevering 16. november}{8}
    \end{ganttchart}
    \caption[Gantt-diagram]{Gantt-diagram for fremdriftsplanen. Løpende rapportskriving pågår gjennom hele prosjektet og er vist med lysere farge.}
    \label{fig:gantt}
\end{figure}

\subsubsection{Risikovurdering og tiltak}
Risiko er vurdert som produktet av sannsynlighet og påvirkning, begge på en skala fra 1 til 5. For personskade er risikoen høyest ved bruk av dreiebenk (12 poeng) og ved sveising (10), og for gjennomføringen at deler blir skadet (12) og at konseptet ikke fungerer (9). De viktigste tiltakene er opplæring og verneutstyr i verkstedet, reservedeler og ferdige printfiler, og at delsystemer testes tidlig. Risikomatrisen, alle hendelsene og tiltakene står i vedlegg~\ref{app:risiko}.
